% Exercise 4.1-4
The simplest change is to leave the algorithm untouched and inspect its result.
Run \textsc{Find-Maximum-Subarray} as before; if the sum it returns is negative,
discard it and return the empty subarray with sum 0 instead, and otherwise
return the subarray unchanged.

This is correct because a maximum subarray with a negative sum can only arise
when every element of $A$ is negative, as shown in exercise 4.1-1, and in that
case the empty subarray, with sum 0, is the better answer under the revised
definition. Whenever $A$ contains at least one nonnegative element, the maximum
sum is itself nonnegative and the empty subarray cannot improve on it. The extra
test costs $\Theta(1)$, so the running time is unchanged.
